\corrigesChapitre{5}{Convexité et inégalités}
\begin{corrige}{1}{Signe de la dérivée seconde}
\begin{questions}
\item La dérivée seconde vaut $2>0$ : convexe sur $\R$.
\item Elle vaut $-2<0$ : concave sur $\R$.
\item Elle vaut $\mathrm e^x>0$ : convexe sur $\R$.
\end{questions}
\end{corrige}
\begin{corrige}{2}{Logarithme et racine}
\begin{questions}
\item $-1/x^2$.
\item Elle est strictement concave sur $]0,+\infty[$ car cette dérivée est négative.
\item $-1/(4x^{3/2})<0$ : la racine est concave sur cet intervalle.
\end{questions}
\end{corrige}
\begin{corrige}{3}{Inflexion}
\begin{questions}
\item $f\,\!^{\prime}=3x^2-3$, donc $f\,\!^{\prime\prime}=6x$.
\item Concave avant $0$, convexe après $0$.
\item La dérivée seconde change de signe en $0$ ; le point est $(0,0)$.
\end{questions}
\end{corrige}
\begin{corrige}{4}{Lire Lorenz}
\begin{questions}
\item Les $40\,\%$ les moins riches possèdent $10\,\%$ de la richesse.
\item $1-0{,}7=0{,}3$, soit $30\,\%$.
\item $0{,}7-0{,}1=0{,}6$, soit $60\,\%$, pour ce groupe de $50\,\%$ de la population.
\end{questions}
\end{corrige}
\begin{corrige}{5}{Une exponentielle amortie}
\begin{questions}
\item $f\,\!^{\prime}=(1-x)\mathrm e^{-x}$, donc $f\,\!^{\prime\prime}=(x-2)\mathrm e^{-x}$.
\item Concave sur $[0,2]$, convexe sur $[2,+\infty[$.
\item $(2,2\mathrm e^{-2})$, car le signe change en $2$.
\end{questions}
\end{corrige}
\begin{corrige}{6}{Une tangente utile}
\begin{questions}
\item $y=1+x$.
\item La convexité donne $\mathrm e^{0{,}2}\geq1{,}2$.
\item La concavité donne $\ln(1{,}2)\leq0{,}2$.
\end{questions}
\end{corrige}
\begin{corrige}{7}{Comparer deux répartitions}
\begin{questions}
\item Les deux valent $0$ en $0$ et $1$ en $1$ ; leurs dérivées $2x$ et $3x^2$ sont positives ou nulles.
\item $L_A-L_B=x^2(1-x)\geq0$.
\item $B$ : pour toute proportion intermédiaire, les personnes les moins riches détiennent une part plus faible.
\end{questions}
\end{corrige}
\begin{corrige}{8}{Déciles}
\begin{questions}
\item $D_9/D_1=3\,600/900=4$.
\item Non. Ce rapport compare des seuils de revenus, pas des masses de revenus cumulées.
\item Il faut les parts cumulées de revenu associées aux proportions cumulées de population, ou les données individuelles ordonnées.
\end{questions}
\end{corrige}
\begin{corrige}{9}{Zéro sans inflexion}
\begin{questions}
\item $f\,\!^{\prime\prime}=12x^2$.
\item Elle s’annule mais ne change pas de signe ; la fonction reste convexe.
\item Rechercher un changement de convexité, notamment un changement de signe de la dérivée seconde, plutôt qu’une simple annulation.
\end{questions}
\end{corrige}
\begin{corrige}{10}{Un modèle admissible ?}
\begin{questions}
\item $L(0)=0,L(1)=1$ et $L\,\!^{\prime}=2-2x\geq0$ sur $[0,1]$.
\item $L-x=x(1-x)\geq0$, strictement à l’intérieur.
\item La courbe est au-dessus de l’égalité et concave. Les groupes les moins riches ne peuvent cumuler plus que leur part de population selon cette définition.
\end{questions}
\end{corrige}
\begin{corrige}{11}{Convexité sans croissance}
\begin{questions}
\item $f\,\!^{\prime}=-\mathrm e^{-x}$ et $f\,\!^{\prime\prime}=\mathrm e^{-x}$.
\item La fonction est strictement décroissante et convexe sur $\R$.
\item Cette fonction constitue un contre-exemple : la convexité porte sur la croissance de la dérivée, pas sur le signe de celle-ci.
\end{questions}
\end{corrige}
\begin{corrige}{12}{Une courbe logistique}
\begin{questions}
\item Le dénominateur et $\mathrm e^{-x}$ sont positifs ; $\mathrm e^{-x}-1$ est positif si $x<0$, négatif si $x>0$.
\item Convexe avant zéro, concave après ; inflexion en $(0,1/2)$.
\item Avant l’inflexion la pente augmente ; après, la croissance ralentit vers une valeur de saturation. Cette interprétation concerne le modèle, pas une loi universelle.
\end{questions}
\end{corrige}
